% !TEX root = IE3_expF_report.tex
\documentclass[lualatex,a4paper,10pt]{jlreq}
% Shared LuaLaTeX preamble.
\usepackage{luatexja-fontspec}
\setmainfont{TeX Gyre Termes}
\setsansfont{TeX Gyre Heros}
\setmonofont{Latin Modern Mono}
\setmainjfont[BoldFont=HaranoAjiGothic-Medium]{HaranoAjiMincho-Regular}
\setsansjfont[BoldFont=HaranoAjiGothic-Bold]{HaranoAjiGothic-Medium}
\usepackage[top=22mm,bottom=22mm,left=23mm,right=23mm]{geometry}
\usepackage{amsmath,amssymb,booktabs,array,longtable,tabularx}
\usepackage{graphicx,xcolor,tikz,circuitikz}
\usetikzlibrary{arrows.meta,positioning,calc}
\usepackage{enumitem,fancyhdr,listings,needspace}
\usepackage[unicode,colorlinks=true,linkcolor=labblue,urlcolor=labteal]{hyperref}
\usepackage{bookmark}
\definecolor{labblue}{RGB}{30,70,112}
\definecolor{labteal}{RGB}{0,100,105}
\definecolor{labgray}{gray}{0.95}
\ModifyHeading{section}{font={\large\sffamily\bfseries\color{labblue}}}
\ModifyHeading{subsection}{font={\normalsize\sffamily\bfseries\color{labteal}}}
\setlength{\parindent}{1\zw}
\setlength{\parskip}{3pt}
\setlength{\emergencystretch}{2em}
\renewcommand{\arraystretch}{1.35}
\setlist{itemsep=3pt,topsep=4pt,leftmargin=2\zw}
\newcolumntype{Y}{>{\raggedright\arraybackslash}X}
\newcolumntype{C}[1]{>{\centering\arraybackslash}p{#1}}
\pagestyle{fancy}
\fancyhf{}
\fancyhead[L]{\small\color{labblue} IE3 後期・電子工学実験}
\fancyhead[R]{\small テーマ\LabCode}
\fancyfoot[C]{\thepage}
\setlength{\headheight}{16pt}
\DeclareOldFontCommand{\rm}{\normalfont\rmfamily}{\mathrm}
\lstset{basicstyle=\small\ttfamily,columns=fullflexible,keepspaces=true,
 breaklines=true,frame=single,backgroundcolor=\color{labgray},
 numbers=left,numberstyle=\tiny,showstringspaces=false,aboveskip=8pt,belowskip=8pt}
\newcommand{\blank}{\rule{22mm}{0.3pt}}
\newcommand{\answerline}{\par\noindent\rule{\linewidth}{0.25pt}\par}
\newcommand{\writing}[1]{\par\Needspace{\dimexpr#1+5mm\relax}\noindent%
 \fcolorbox{labblue!35}{white}{\parbox[c][#1][t]{\dimexpr\linewidth-2\fboxsep-2\fboxrule\relax}{\vspace{2mm}\noindent\strut\hfill\par}}\par}
\newcommand{\hint}[1]{\par\smallskip\noindent{\sffamily\bfseries\color{labteal}記述の視点：}#1\par}
\newcommand{\note}[1]{\par\smallskip\noindent\fcolorbox{labblue!45}{labblue!5}{%
 \parbox{\dimexpr\linewidth-2\fboxsep-2\fboxrule\relax}{#1}}\par\smallskip}
\newcommand{\eqerror}{\varepsilon=\frac{x_{\mathrm{meas}}-x_{\mathrm{th}}}{x_{\mathrm{th}}}\times100\ \%}
\newcommand{\LabSource}{徳山工業高等専門学校 情報電子工学科，
 『2026年度 電子工学実験（3年後期）テキスト』，テーマ\LabCode，
 \expandafter\path\expandafter{\SourcePdf}。}
\newcommand{\reporttitle}{
 \noindent{\small\sffamily\color{labteal}実験報告書・記入ガイド}\par
 \noindent{\LARGE\sffamily\bfseries \LabCode\quad\LabTitle}\par\smallskip
 \noindent 氏名：\blank\quad 班：\blank\quad 実験日：\blank\par
 \note{目的・原理・手順を確認し、実測値と自分の考察を記入する。
 考察のガイドは、検討する事象・使う測定結果・記述の方針を示す。}
}
\newcommand{\prelabtitle}{
 \noindent{\small\sffamily\color{labteal}事前学習・前提知識プリント}\par
 \noindent{\LARGE\sffamily\bfseries \LabCode\quad\LabTitle}\par\smallskip
 \noindent 氏名：\blank\quad 班：\blank\quad 予習日：\blank\par
 \note{用語、回路の動き、計算、測定表への記入の順に読む。
 例題の値は説明用であり、実験結果ではない。}
}
\newcommand{\tocrow}[3]{\hyperref[#1]{#2}& #3 &\pageref*{#1}\\[2mm]}
\newcommand{\documentmap}[1]{
 \section*{目次と概要}
 \noindent 青色の項目をクリックすると該当箇所へ移動する。\par
 {\small\begin{longtable}{@{}p{0.29\linewidth}p{0.61\linewidth}r@{}}
 \toprule {\color{labblue}\bfseries 項目} & {\color{labblue}\bfseries 読むと分かること} & 頁\\\midrule
 \endhead #1\bottomrule
 \end{longtable}}
 \clearpage
}
\newenvironment{terms}{\par\smallskip\begin{longtable}{@{}p{0.23\linewidth}p{0.73\linewidth}@{}}
 \toprule {\color{labteal}\bfseries 用語}&{\color{labteal}\bfseries 意味・回路での役割}\\\midrule\endhead}
 {\bottomrule\end{longtable}}
\newcommand{\term}[2]{\textbf{#1}&#2\\[2mm]}
\newcommand{\guidepoint}[4]{
 \Needspace{32mm}\subsection{#1}
 \noindent{\bfseries\color{labteal}考える事象：}#2\par
 \noindent{\bfseries\color{labteal}参照する結果：}#3\par
 \noindent{\bfseries\color{labteal}記述の方針：}#4\par
}
\newcommand{\reportclosing}[1]{
 \clearpage\section{結論のまとめ方}\label{sec:closing}
 \hint{#1}
 \ConclusionGuide
 \begin{enumerate}
 \item 目的に対応する最も重要な実測結果を、測定条件と数値を添えて述べる。
 \item 原理と一致した点・一致しなかった点を示す。原因は確認済みの事実と推測を分ける。
 \item 実施範囲で言えることと、未確認の条件を一文ずつまとめる。
 \end{enumerate}
 \note{書き出しの型：「○○の条件で△△を測定したところ、□□となった。
 これは◇◇と整合する。ただし××は確認しておらず、一般化には追加測定が必要である。」
 空欄は自分の実測値と根拠で埋める。}
 \noindent 結論（実験後に記入）：\writing{30mm}
 \noindent 改善案・次に確認したい条件：\writing{16mm}
 \section*{参考文献}\label{sec:references}
 \begin{enumerate}
 \item \LabSource
 \item 使用したデータシート・書籍・ウェブ資料（著者、題名、該当箇所、URL、参照日）を追加する。
 \end{enumerate}
}
\newcommand{\prelabclosing}{
 \section*{準備・出典}\label{sec:prelabsource}
 \begin{itemize}[nosep]
 \item 資料、筆記具、記録用紙、電卓と、実験条件・機器設定の記録欄。
 \item 確認問題の解答と、分からない用語・回路点への具体的な質問。
 \end{itemize}
 {\small\noindent\LabSource\par}
}
\newcommand{\simpleflow}[1]{
 \begin{center}
 \begin{tikzpicture}[node distance=5mm and 5mm,
 every node/.style={font=\small},box/.style={draw=labblue,fill=labblue!4,rounded corners=1mm,
 align=center,minimum height=10mm,text width=30mm}]
 #1
 \end{tikzpicture}
 \end{center}
}

\newcommand{\LabCode}{F}
\newcommand{\LabTitle}{カウンタ回路}
\newcommand{\SourcePdf}{IE3_expF_A4.pdf}
\newcommand{\ConclusionGuide}{各段の循環と主要な境界動作を根拠に、24時間時計として確認できた範囲を述べる。
基準周期の測定値を添え、「論理が正しいこと」と「実時間で正確なこと」を別々にまとめる。\par}
\begin{document}
\reporttitle
\documentmap{
\tocrow{sec:auto1}{1 目的}{分周と桁上げで構成する24時間時計を実験で確かめる。}
\tocrow{sec:auto2}{2 原理}{分周と桁上げで構成する24時間時計の仕組みと成立条件を整理する。}
\tocrow{sec:auto3}{3 装置・方法}{機器・定数・測定点を確認し、資料の順序で操作する。}
\tocrow{sec:auto4}{4 結果}{各桁の計数、境界動作、クロック周期を表や図に記録する。}
\tocrow{sec:auto5}{5 考察：結果を根拠に説明する}{各桁の計数、境界動作、クロック周期を根拠に、比較と原因の検討を進める。}
\tocrow{sec:closing}{6 結論のまとめ方}{主要な結果、成立範囲、未確認の条件を短くまとめる。}
\tocrow{sec:references}{参考文献}{使用した資料と外部の根拠を示す。}
}

\section{目的}\label{sec:auto1}
カウンタ、デコーダ、7セグメント表示器と555のクロックを用いて24時間時計を構成し、分周・桁上げ・リセットの仕組みを理解する。
\section{原理}\label{sec:auto2}
時計の各桁は、秒・分の一の位が10進、十の位が6進、時が00〜23の循環となる。
7490は内部の2分周と5分周を接続して10進にする。7492は2分周と6分周を持ち、全体を接続すると12分周であるため、6進に使う部分と入力端子を資料で確認する。
各カウンタの規定のエッジで計数し、桁上げ信号を次段へ送る。
時の一の位4と十の位2を検出して24を零に戻す。
\[
 f_{\rm clock}\simeq\frac1{0.693(R_A+2R_B)C}
\]
は555非安定回路の近似である。
\simpleflow{
\node[box] (a) {555\\約1 Hz};
\node[box,right=of a] (b) {秒\\10進・6進};
\node[box,right=of b] (c) {分\\10進・6進};
\node[box,right=of c] (d) {時\\24進};
\draw[-{Stealth}] (a)--(b);\draw[-{Stealth}] (b)--(c);\draw[-{Stealth}] (c)--(d);
}
\section{装置・方法}\label{sec:auto3}
\begin{enumerate}
\item 555クロックと手動リセットを先に確認する。電源、ICの方向、リセット入力の初期状態を確認する。
\item 一桁ずつカウンタ、デコーダ、表示器を配線し、0〜9または0〜5の循環を確認する。
\item 資料指定の桁上げ端子を接続する。10進は$Q_D$、6進は$Q_C$の波形と次段の計数エッジを確認する。
\item 時の24検出と手動リセットを接続し、全桁が零になることを確認する。
\item 通常クロックで動作させた後、指定の加速試験により分・時・24時の境界を確認する。
\end{enumerate}
実際のIC型番・回路図・555のR、C・クロック周期：
\writing{22mm}
\clearpage
\section{結果}\label{sec:auto4}
表1：各段の分周確認。
\par\noindent\begin{tabularx}{\linewidth}{YYYY}
\toprule 段 & 使用IC・入力 & 表示範囲 & 桁上げ波形・周期\\\midrule
秒一の位&&&\\秒十の位&&&\\分一の位&&&\\分十の位&&&\\時&&&\\\bottomrule
\end{tabularx}
表2：境界動作。予想値と観測値を区別する。
\par\noindent\begin{tabularx}{\linewidth}{YYYY}
\toprule 試験直前 & 予想される次表示 & 観測した表示 & 試験条件\\\midrule
00:00:09 &00:00:10&&\\00:00:59&00:01:00&&\\
00:59:59&01:00:00&&\\23:59:59&00:00:00&&\\
手動リセット&00:00:00&&\\\bottomrule
\end{tabularx}
図1：7490、7492の各出力と桁上げのタイミング図。資料の計数順序を確認する。
\writing{45mm}
図2：24検出、手動リセット、各桁の接続を含む回路図（別紙追加可）。
\writing{40mm}
\clearpage\section{考察：結果を根拠に説明する}\label{sec:auto5}
\guidepoint{各段の分周と桁上げ}
{10進・6進の循環と、次段を進める信号が正しいか。}
{表1の各段、7490/7492のタイミング図、入力と桁上げの周期。}
{入力端子・桁上げ端子・計数エッジを示し、周期比を求める。7492の全体12分周と使用部分6分周を区別する。}
\guidepoint{24時と手動リセット}
{23:59:59の次が00:00:00となり、リセット後に再び進むか。}
{表2の境界試験、24検出回路、リセット波形、再開時の表示。}
{各桁に伝わる桁上げと時のリセットを追う。短い検出パルス・遅延の影響を候補にし、表示だけでパルス幅を推定しない。}
\guidepoint{クロックの設計と精度}
{実測周期がRC式と約1秒の目標に一致するか。}
{実際のR、C、555の波形・周期、計時の測定時間。}
{理論と実測を比較し、周期ずれの累積を見積もる。加速境界試験と実時間精度の評価を分け、長時間測っていない場合は精度の結論を限定する。}
\guidepoint{段階的な組立てと異常の切り分け}
{表示・桁上げの異常がどの段から起こったか。}
{単独段の確認記録、配線修正前後の表示、ICとデコーダの接続図。}
{症状→確認した端子→変更→再確認の順に書く。動いた最終状態だけでなく、原因を特定できた根拠を残す。}
\subsection{資料の課題・追加検討}
\begin{enumerate}
\item 各カウンタの入力と出力の周波数関係を、タイミング図で説明する。
\item 24でリセットする論理と、短いリセットパルス・伝搬遅延の影響を説明する。
\item 実測のクロック周波数とR、Cによる計算値を比較し、約1 Hzになる抵抗を設計する。
\item 作業で発生した誤配線や表示異常、その発見方法と改善案を記す。
\end{enumerate}
\writing{25mm}
\reportclosing{時計の分周比と24時間動作の確認結果をまとめる。}
\end{document}
